% Derived from research/information/theory.tex; original preserved.
% Classical information/risk toolbox and explicit boundary examples.
% This is a section body, not a new-novelty claim.
% Parent requirements: amsmath, amssymb, amsthm; proposition and remark environments.
% Bibliography entries are supplied in information_references.bib.
\section{What information gained by a representation does and does not buy}
\label{sec:p7-information-boundaries}

All logarithms are natural. Representations and the target are random
variables on standard Borel spaces; conditional laws exist. The target is
finite except where real regression is stated explicitly. Population risks
are evaluated on a fresh example. For learned encoders, all statements can
be applied conditional on a fixed training history, provided the fresh
example has the stipulated conditional law. Conditioning does not assert
that the encoder is independent of its training examples.

Write $B_\ell(Z)=\inf_g\mathbb E\ell(Y,g(Z))$, where the infimum
runs over all measurable decoders with the appropriate action space.
A deterministic refinement has $Z=\pi(Z')$ almost surely. More generally,
the conclusions stated for refinements hold when $Y-Z'-Z$ is a Markov
chain under the specified coupling. Appending a feature is a refinement;
retraining and replacing an entire hidden layer need not be one.
The results in this section are classical identities, elementary
consequences, and explicit counterexamples. Proper-loss information and
Bregman interpretations are established in
\cite[Theorems 8 and 10]{p7reid2011} and
\cite[Theorem 1 and Section 2.2]{p7gneiting2007}.

\begin{proposition}[Information gain includes a possible forgetting term]
\label{prop:p7-log-forgetting}
For finite $Y$, logarithmic loss gives $B_{\log}(Z)=H(Y\mid Z)$.
For any pair of representations on a common probability space,
\[
 B_{\log}(Z)-B_{\log}(Z')
 =I(Y;Z'\mid Z)-I(Y;Z\mid Z').
\]
For a refinement, the second term vanishes, and the Bayes improvement is
exactly $I(Y;Z'\mid Z)$.
\end{proposition}
\begin{proof}
Conditional expected log loss at $Z=z$ equals the entropy of the true
conditional distribution plus its KL divergence to the prediction.
The divergence is nonnegative and is zero at the true distribution.
Next expand both conditional mutual informations, using
$H(Y\mid Z,Z')$ as their common subtracted term. Finally the Markov
condition makes $I(Y;Z\mid Z')=0$.
\end{proof}

Thus positive conditional mutual information concerns an optimal
unrestricted decoder; it is not a decrease promised for a trained decoder.
Nor does adding and then reoptimizing neurons automatically preserve the
information-refinement condition.

\begin{proposition}[Squared and classification risk under refinement]
\label{prop:p7-basic-bounds}
Suppose $Y-Z'-Z$ is Markov. If $Y$ is real with
$\mathbb E Y^2<\infty$, put $m=\mathbb E[Y\mid Z]$ and
$m'=\mathbb E[Y\mid Z']$. Then
\[
 \Delta_{\mathrm{sq}}:=B_{\mathrm{sq}}(Z)-B_{\mathrm{sq}}(Z')
 =\mathbb E(m'-m)^2.
\]
If additionally $Y\in[a,b]$ almost surely and
$J=I(Y;Z'\mid Z)$, then
\[
 0\le\Delta_{\mathrm{sq}}\le\tfrac12(b-a)^2 J.
\]
For finite-label classification with zero--one loss,
\[
 0\le\Delta_{01}:=B_{01}(Z)-B_{01}(Z')\le\sqrt{J/2}.
\]
For $Y\in\{0,1\}$, both constants are asymptotically sharp among
refinements. Here the Brier loss means $(Y-p)^2$, not the sum over both
class coordinates. If the binary coarse posterior satisfies
$m\in[\alpha,1-\alpha]$ almost surely, $0<\alpha\le1/2$, then also
\[
 J\le\frac{\Delta_{\mathrm{sq}}}{\alpha(1-\alpha)}.
\]
The last bound is a convenient reverse bound, not a sharp constant claim.
\end{proposition}
\begin{proof}
Conditional expectation minimizes squared loss. The tower identity and
the Markov condition give $\mathbb E[m'\mid Z]=m$, hence
$\mathbb E[m'm]=\mathbb E m^2$. Subtracting the two Bayes risks gives
$\mathbb E(m')^2-\mathbb E m^2=\mathbb E(m'-m)^2$.

Write $P'=P_{Y\mid Z'}$ and $P=P_{Y\mid Z}$. Under the Markov
condition, $J=\mathbb E\mathrm{KL}(P'\Vert P)$.
The mean difference for a variable in $[a,b]$ is at most
$(b-a)\mathrm{TV}(P',P)$. Pinsker's inequality, in nats, states
$\mathrm{TV}(P',P)^2\le\mathrm{KL}(P'\Vert P)/2$.
Integration proves the squared-risk bound.
The conditional best classification accuracy is $\max_y P(y)$.
Since $\max_y P'(y)-\max_y P(y)\le\mathrm{TV}(P',P)$,
integration, Cauchy--Schwarz, and Pinsker give the upper bound.
The nonnegative sign follows from conditional Jensen, since
$P=\mathbb E[P'\mid Z]$ and the maximum is convex.

For sharpness let $Z$ be constant and let $m'=1/2\pm\epsilon$ with
equal probabilities. Then $\Delta_{\mathrm{sq}}=\epsilon^2$,
$\Delta_{01}=\epsilon$, and
$J=\log2-h(1/2+\epsilon)=2\epsilon^2+O(\epsilon^4)$, where
$h(p)=-p\log p-(1-p)\log(1-p)$.
Finally $\log u\le u-1$ gives
$\operatorname{kl}(q,p)\le(q-p)^2/[p(1-p)]$ for $0<p<1$.
Apply this conditionally and integrate.
\end{proof}

The bounded-response implication is an instance of the classical
information-to-conditional-mean inequality; see the earlier absolute-mean
version in \cite[Lemma 4.4]{p7tao2006} and the explicit sharp squared-mean formulation in
\cite[Corollary 7.11]{p7polyanskiy2024}. Pinsker and reverse-divergence
comparisons are established tools \cite[Eq.~(1), Section VI]{p7sason2016}.
These statements are not new neuron-growth guarantees.

\begin{proposition}[Function preservation and representation preservation differ]
\label{prop:p7-birth}
If $Z'=F(Z)$ and $Z=G(Z')$ almost surely for measurable $F,G$, then
the two representations generate the same sigma field and
$I(Y;Z')=I(Y;Z)$, $I(Y;Z'\mid Z)=0$.
Exact neuron duplication before symmetry-breaking perturbations has
this property for the complete hidden representation.

Nevertheless, a function-preserving neural birth can have strictly
positive representation information gain. Fix $0<\delta<1/2$,
let $X$ be uniform on $\{-1,1\}$, and set
\[
 \mathbb P(Y=1\mid X=1)=1-\delta,\qquad
 \mathbb P(Y=1\mid X=-1)=\delta.
\]
Start with constant representation $Z=1$ and prediction $f=1/2$.
Append the ReLU feature $r(X)=\max\{X,0\}$ with output coefficient
zero. The new representation is $Z'=(1,r(X))$, while the new prediction
is still $f'=1/2+0r(X)=1/2$. Then
\[
 I(Y;Z'\mid Z)=\log2-h(\delta)>0,
 \qquad \Delta_{\mathrm{sq}}=(1/2-\delta)^2.
\]
The actual prediction risk at birth is unchanged for every loss for
which it is defined. In particular its squared risk is $1/4$, its log
risk is $\log2$, and its classification error under the same fixed
tie rule is $1/2$, both before and after birth.
\end{proposition}
\begin{proof}
Mutual measurability gives equality of sigma fields, hence equality of
all conditional laws. In the neural example $r(X)$ determines $X$.
The marginal label is balanced, with entropy $\log2$, and its
conditional entropy given $X$ is $h(\delta)$. Its conditional mean is
$\delta+(1-2\delta)r(X)$, at distance $1/2-\delta$ from $1/2$.
Propositions~\ref{prop:p7-log-forgetting}--\ref{prop:p7-basic-bounds}
give the asserted Bayes gains. The realized predictions coincide
pointwise, which proves equality of their realized risks directly.
\end{proof}

Function-preserving duplication is established in Net2Net
\cite[Section 2.3]{p7chen2016}; the distinction here is between the
unchanged output function and what an unrestricted decoder could recover
from the enlarged hidden vector. In the displayed example even an affine
decoder can express the new Bayes predictor, but the birth operation has
not fitted its coefficients.

\begin{proposition}[Four explicit failures of information-only criteria]
\label{prop:p7-counterexamples}
The following claims hold, each under a deterministic refinement.
\begin{enumerate}
\item Activation entropy can be arbitrarily large while task information
and every Bayes-risk improvement are zero. Let $U$ be uniform on
$\{1,\ldots,M\}$, independent of $Y\sim\mathrm{Ber}(1/2)$.
Take $Z=1$ and $Z'=\mathrm{ReLU}(U)=U$. Then
$H(Z')=\log M$, $I(Y;Z'\mid Z)=0$, and the conditional target law
is unchanged.
\item Positive label information need not improve classification. Let
$Z=1$, and let $Z'$ be equiprobable with binary label posteriors
$1/10$ and $3/10$. Strict concavity gives
$J=h(1/5)-\{h(1/10)+h(3/10)\}/2>0$, but both Bayes classification
errors are $1/5$, since the Bayes action stays zero. This rules out
a universally positive classification-gain lower bound at every
positive information level; it does not rule out all bounds at larger
information levels.
\item Without restricting the response or the decision loss, even a
fixed positive information gain need not reduce squared Bayes risk.
Let $U\sim\mathrm{Ber}(1/2)$ and independent $S$ be uniform on
$\{-1,1\}$. Set $Y=S(1+U)$, $Z=1$, and $Z'=U$.
Then $U=|Y|-1$, so $J=\log2$, but both conditional means are zero
and both squared Bayes risks equal $\mathbb E Y^2=5/2$.
\item Fixed structural information cost does not imply useful risk
reduction. For independent equiprobable $S\in\{-1,1\}$ and uniform
$U\in\{1,\ldots,M\}$, let $X=(S,U)$, $Z=1$, $Z'=X$, and
$Y\mid X\sim\mathrm{Ber}(1/2+\epsilon S)$ with
$0<\epsilon<1/2$. Then
\[
 I(X;Z'\mid Z)=\log(2M),\quad
 J=\log2-h(1/2+\epsilon),\quad
 \Delta_{\mathrm{sq}}=\epsilon^2,\quad\Delta_{01}=\epsilon.
\]
All three Bayes gains tend to zero as $\epsilon\downarrow0$ while
the structural information cost is fixed and can be arbitrarily large.
\end{enumerate}
\end{proposition}
\begin{proof}
Independence proves the first item. The second follows by averaging the
two conditional errors and entropies. The third follows from independence
and the zero mean of $S$. In the fourth, $Z'$ reveals $X$, so its input
information is $H(X)=\log(2M)$. The label is balanced, the conditional
entropy is $h(1/2+\epsilon)$, the posterior mean is $1/2+\epsilon S$,
and the fine classification error is $1/2-\epsilon$. These observations
give all displayed formulas.
\end{proof}

Even for binary labels there is no distribution-free positive constant
$c$ with $\Delta_{\mathrm{sq}}\ge cJ$ for all refinements.
Indeed let $X=Y\sim\mathrm{Ber}(p)$, $Z=1$, $Z'=X$. Then
$\Delta_{\mathrm{sq}}=p(1-p)$ and $J=h(p)$, whose ratio tends to
zero as $p\downarrow0$. This is only a failure of a uniform linear
comparison; it is not a claim that no nonlinear comparison exists.

\begin{proposition}[The exact gap between Bayes value and fitted risk]
\label{prop:p7-fitting-gaps}
For $t\in\{0,1\}$, fix representation $Z_t$, an admissible decoder
class $\mathcal G_t$, and its fitted member $g_t$. Assume all risks and
infima below are finite. Let $R_t(g)=\mathbb E\ell(Y,g(Z_t))$,
$B_t=B_\ell(Z_t)$, and define
\[
 A_t=\inf_{g\in\mathcal G_t}R_t(g)-B_t\ge0,\qquad
 F_t=R_t(g_t)-\inf_{g\in\mathcal G_t}R_t(g)\ge0.
\]
Then
\[
 R_0(g_0)-R_1(g_1)
 =(B_0-B_1)+(A_0-A_1)+(F_0-F_1).
\]
If $\widehat R_t$ is an empirical criterion, an exact further split is
\[
 F_t=O_t+E_t,
 \quad O_t=\widehat R_t(g_t)-\inf_{g\in\mathcal G_t}\widehat R_t(g)\ge0,
\]
\[
 E_t=R_t(g_t)-\widehat R_t(g_t)
       +\inf_{g\in\mathcal G_t}\widehat R_t(g)
       -\inf_{g\in\mathcal G_t}R_t(g).
\]
Here $E_t$ can have either sign. If
$U_t=\sup_{g\in\mathcal G_t}|R_t(g)-\widehat R_t(g)|<\infty$,
then $|E_t|\le2U_t$.
\end{proposition}
\begin{proof}
Add and subtract each infimum to obtain the identities. Uniform
deviation bounds both the fitted function's risk difference and the
difference of the two infima by $U_t$ in absolute value.
\end{proof}

$A_t$ is representation-decoder approximation error; $F_t$ is the
realized population fitting gap, not a pure optimization error.
Only $O_t$ is the empirical optimization residual in the displayed split.
The algebra holds for data-dependent encoders and classes after they are
fixed, but obtaining a small probabilistic bound for $U_t$ then requires
appropriate uniform control or fresh data. The identities do not supply
one. A larger realizable function class can lower its infimum while
finite optimization fails to lower the fitted risk.

\begin{proposition}[A total information budget requires nested retention]
\label{prop:p7-chain-budget}
Let $Z_0,\ldots,Z_T$ be deterministic successive refinements of $X$.
Then
\[
 \sum_{t=1}^T I(Y;Z_t\mid Z_{t-1})
 = I(Y;Z_T\mid Z_0)
 = B_{\log}(Z_0)-B_{\log}(Z_T)
 \le I(Y;X\mid Z_0).
\]
If all information quantities are finite and
$C=\sum_{t=1}^T I(X;Z_t\mid Z_{t-1})$, then also
\[
 C=I(X;Z_T\mid Z_0),\qquad
 B_{\log}(Z_0)-B_{\log}(Z_T)\le C.
\]
Write $M=\min\{I(Y;X\mid Z_0),C\}$. For bounded real
$Y\in[a,b]$, the endpoint squared Bayes improvement is at most
$(b-a)^2M/2$. The corresponding finite-label classification bound is
$\sqrt{M/2}$.
\end{proposition}
\begin{proof}
Each earlier representation is a function of every later one. Apply the
mutual-information chain rule and telescope the conditional entropies
of the finite target. Apply the same chain rule to $X$ (via mutual
information, without requiring finite differential entropies).
Conditional data processing gives
$I(Y;Z_T\mid Z_0)\le I(Y;X\mid Z_0)$ and
$I(Y;Z_T\mid Z_0)\le I(X;Z_T\mid Z_0)$.
Proposition~\ref{prop:p7-basic-bounds} applied once to the endpoint pair
proves the remaining bounds.
\end{proof}

The same chain-rule conclusions hold for an explicitly specified joint
Markov degradation chain $Y-X-Z_T-\cdots-Z_0$. They need not hold for
arbitrary successive fitted representations. For example, take a balanced
binary $Y=X$ and alternate $Z_t=X$ at odd $t$ and a constant at even $t$,
starting with constant $Z_0$. Every odd transition contributes $\log2$
to $I(Y;Z_t\mid Z_{t-1})$; every even transition contributes zero.
Their sum grows without bound as the number of cycles grows, although
the final representation after any full cycle has no label information.
The omitted reverse conditional information accounts exactly for each
forgetting step. A cumulative transcript restores nesting, but describes
the information of that retained transcript, not automatically that of
the deployed network. None of these chain-rule facts controls optimizer
descent or charges finite-precision storage, feature-search computation,
or parameter updates.

For a learned sequence, a valid conditional version fixes one common
training history $H$ containing the randomness that specifies every
encoder, and uses a fresh evaluation sample independent of $H$.
Apply the preceding propositions under each conditional law given $H$,
then average. Conditioning separately on a growing $H_t$ at each stage
without accounting for the newly learned information is not the same
argument. In particular the randomness specifying the next encoder
must be conditioned on, or represented explicitly as part of the
observation, before applying a fixed-encoder Bayes comparison.
